
\newpage
\section{Метод расчета параметров влияния волн на берег}
\label{part:sea_stuff}

\epigraph{
    Change my pitch up!\\
    Smack my \cancel{bitch} beach up!
}{\textit{--~The Prodigy}}

Описывая влияние моря на бреговую линию его можно разделить на две группы факторов: механических, определяемых влиянием динамических нагрузок формируемых морскими волнами, а также химических, вызываемых агрессивным воздействием минерального состава морской воды.
Влияние химических факторов проявляется в последовательной снижении физико-механических характеристик прибрежных конструкций и заложено в срок эксплуатации объекта на этапе его проектирования.
Исходя их этого, в последующей части работы будут рассмотрены только механические факторы.

Учитывая то, что приоритет в выполняемой работе был отдан механизмам уязвимости береговой линии, то наиболее удачный показатель, описывающий потенциал к ее разрушению должен определять энергию, передаваемую морскими волнами объектам при их столкновении.

Численное выражение этого показателя естественным образом зависит от интенсивности, и направления волновых процессов приходящихся на рассматриваемую точку береговой линии.
Математическое описание и алгоритмизация этих процессов не возможны без раскрытия базовой теории формирования морских волн и описания механизмов их трансформации на подходе к береговой линии, чему посвящен следующий подраздел работы.   
